\chapter{Light-Matter Interaction}
\label{ch:method}

\section{Theory of Light--Matter Interaction}
\subsection{Fermi's golden rule}

In this section, we examine the interaction between light and matter using \textbf{Fermi's golden rule}. The transition rate for an optical excitation from the ground state $|GS\rangle$, with fully occupied valence bands, to an exciton state $|\Psi_{S, \boldsymbol{Q}}^{X}\rangle$ with band index $S$ and center-of-mass momentum $\boldsymbol{Q}$ is given by
\begin{equation}
\Gamma=\frac{2\pi}{\hslash}
\left|
\langle\Psi_{S, \boldsymbol{Q}}^{X}|\hat{H}_I|GS\rangle
\right|^2
\delta(E^X_{S, \boldsymbol{Q}}-\hslash\omega).
\label{Fermi}
\end{equation}
The delta function enforces energy conservation between the exciton energy $E^X_{S, \boldsymbol{Q}}$ and the photon energy $\hslash\omega$. Within the pseudospin model~\cite{simbulan2021selective, peng2022tailoring, sanchez2024enhanced}, the exciton state is written as
\begin{equation}
|\Psi_{S, \boldsymbol{Q}}^{X}\rangle
=
\frac{1}{\sqrt{\Omega}} 
\sum_{vc\boldsymbol{k}} 
\Lambda_{S,\boldsymbol{Q}}(vc\boldsymbol{k})
\hat{c}^\dagger_{c,\boldsymbol{k}+\boldsymbol{Q}} 
\hat{h}^\dagger_{v,-\boldsymbol{k}} 
|GS\rangle.
\label{EXstate}
\end{equation}
Here, the particle operators $\hat{c}_{c,\boldsymbol{k}}^\dagger$ and $\hat{h}_{v,-\boldsymbol{k}}^\dagger$ create an electron in the conduction band with crystal momentum $\boldsymbol{k}$ and a hole in the valence band with crystal momentum $-\boldsymbol{k}$, respectively. The coefficient $\Lambda_{S,\boldsymbol{Q}}(vc\boldsymbol{k})$ represents the amplitude of the electron--hole configuration for exciton band $S$, and $\Omega$ is the area of the TMD monolayer.
The matrix element in Eq.~(\ref{Fermi}) can be expressed, within the \textbf{rotating wave approximation}~\cite{griffiths2018introduction}, as
\begin{equation}
\langle\Psi_{S, \boldsymbol{Q}}^{X}|\hat{H}_I|GS\rangle 
=
\frac{1}{\sqrt{\Omega}} 
\sum_{vc\boldsymbol{k}} 
\Lambda^*_{S,\boldsymbol{Q}}(vc\boldsymbol{k}) 
\langle\psi_{c, \boldsymbol{k}+\boldsymbol{Q}}|H_I|\psi_{v, \boldsymbol{k}}\rangle,
\label{eq:matrix_element}
\end{equation}
where, for weak fields, the light--matter interaction Hamiltonian is given by~\cite{peng2022tailoring, sanchez2024enhanced}
\[
\hat{H}_I=\frac{|e|}{2m_e}\boldsymbol{A}(\boldsymbol{r}, t)\cdot\boldsymbol{p},
\]
with $\boldsymbol{A}(\boldsymbol{r}, t)$ the vector potential in the \textbf{Coulomb gauge}, $\boldsymbol{p}$ the momentum operator, and $m_e$ ($e$) the electron mass (charge).
If the vector potential is expanded in a \textbf{plane-wave basis} as in Eq.~(\ref{planwaveinC}), the matrix element in Eq.~(\ref{eq:matrix_element}) becomes
\begin{align}
\langle\psi_{c, \boldsymbol{k}+\boldsymbol{Q}}|H_I|\psi_{v, \boldsymbol{k}}\rangle
&= \frac{|e|}{2m_e} 
\sum_{\boldsymbol{q}_\parallel} 
\mathcal{A}(\boldsymbol{q}_\parallel) 
\left[
\hat{\boldsymbol{\varepsilon}} \cdot 
\langle\psi_{c, \boldsymbol{k}+\boldsymbol{Q}}|
e^{i\boldsymbol{q}_\parallel \cdot \boldsymbol{r}} 
\hat{\boldsymbol{p}} 
|\psi_{v, \boldsymbol{k}}\rangle
\right] e^{i\omega t} \nonumber \\
&\approx \frac{|e|}{2m_e} 
\sum_{\boldsymbol{q}_\parallel} 
\mathcal{A}(\boldsymbol{q}_\parallel) 
\left[
\hat{\boldsymbol{\varepsilon}} \cdot 
\delta_{\boldsymbol{q}_\parallel,\boldsymbol{Q}} 
\langle\psi_{c, \boldsymbol{k}}| 
\hat{\boldsymbol{p}} 
|\psi_{v, \boldsymbol{k}}\rangle
\right] e^{i\omega t},
\end{align}
where the approximation follows from the \textbf{electric dipole approximation}. In this approximation, the Bloch state wave functions are written as $\psi_{n, \boldsymbol{k}}(\boldsymbol{r})=u_{n, \boldsymbol{k}}(\boldsymbol{r}) e^{i\boldsymbol{k}\cdot \boldsymbol{r}}$, and the plane-wave components combine to yield the Kronecker delta $\delta_{\boldsymbol{q}_\parallel,\boldsymbol{Q}}$, enforcing the equality between the in-plane wavevector and the exciton center-of-mass momentum.

Substituting this result back, the matrix element in Eq.~(\ref{eq:matrix_element})~\cite{peng2022tailoring, sanchez2024enhanced} becomes
\begin{equation}
\langle\Psi_{S, \boldsymbol{Q}}^{X}|\hat{H}_I|GS\rangle 
=
\frac{E_g}{2i\hslash}
\mathcal{A}(\boldsymbol{Q}) 
\left(
\hat{\boldsymbol{\varepsilon}} \cdot 
\boldsymbol{D}^{X*}_{S, \boldsymbol{Q}}
\right) e^{i\omega t},
\label{firstmatrixelement}
\end{equation}
where $\boldsymbol{D}^{X}_{S, \boldsymbol{Q}}$ is the exciton transition dipole moment, superscript $^*$ denotes complex conjugation, and $E_g$ is the band gap at the $K/K'$ valleys. This expression shows that optical excitation is governed by the alignment between the light polarization and the transition dipole moment. The transition rate is therefore proportional to both the magnitude of the transition dipole moment and the magnitude of the vector potential. Under the excitaion of VVBs, the optical matrix element is defined as
\begin{align}
\tilde{M}^{\sigma, \ell, \sigma', \ell'}_{S, \boldsymbol{Q}}
&\equiv
\frac{1}{\sqrt{\Omega}} 
\sum_{vc\boldsymbol{k}} 
\Lambda^*_{S,\boldsymbol{Q}}(vc\boldsymbol{k}) 
\left\langle
\psi_{c, \boldsymbol{k}+\boldsymbol{Q}}
\left|
\frac{|e|}{2m_e}
\boldsymbol{\mathcal{A}}^{\sigma, \ell, \sigma', \ell'}(\boldsymbol{Q};\alpha,\beta)
\cdot
\boldsymbol{p}
\right|
\psi_{v, \boldsymbol{k}}
\right\rangle \nonumber \\
&\approx
\frac{E_g}{2i\hslash}
\mathcal{A}^{\sigma, \ell, \sigma', \ell'}(\boldsymbol{Q};\alpha,\beta)
\left(
\hat{\boldsymbol{\varepsilon}} \cdot 
\boldsymbol{D}^{X*}_{S, \boldsymbol{Q}}
\right),
\label{matrixelement}
\end{align}
where the approximation follows from the electric dipole approximation. The corresponding Fermi's golden rule can then be written as~\cite{sanchez2024enhanced}
\begin{equation}
\Gamma^{\sigma, \ell, \sigma', \ell'}_{S, \boldsymbol{Q}}
=
\frac{2\pi}{\hslash}
\left|
\tilde{M}^{\sigma, \ell, \sigma', \ell'}_{S, \boldsymbol{Q}}
\right|^2
\rho(\hslash\omega=E^X_{S, \boldsymbol{Q}}) \\  
\propto
\left|
\mathcal{A}^{\sigma, \ell, \sigma', \ell'}(\boldsymbol{Q};\alpha,\beta)
\left(\hat{\boldsymbol{\varepsilon}} \cdot 
\boldsymbol{D}^{X}_{S, \boldsymbol{Q}}
\right)\right|^2. \label{transitionrate}
\end{equation}
The transition dipole moment, obtained from density functional theory  combined with the Bethe--Salpeter equation, is given in the small-$\boldsymbol{Q}$ limit by Ref~\cite{peng2022tailoring, sanchez2024enhanced}:
\begin{figure}[h] 
\centering
\begin{minipage}{0.45\textwidth}
    \centering
    \textbf{Valley excitons (bright)}
    \begin{equation}
    \begin{aligned}
        \boldsymbol{D}^{X}_{K, \boldsymbol{Q}} &= \frac{D^X_B}{\sqrt{2}} (\hat{\boldsymbol{x}}+i\hat{\boldsymbol{y}}), \\
        \boldsymbol{D}^{X}_{K', \boldsymbol{Q}} &= \frac{D^X_B}{\sqrt{2}} (\hat{\boldsymbol{x}}-i\hat{\boldsymbol{y}})
    \end{aligned}
    \end{equation}
\end{minipage}
\hfill
\begin{minipage}{0.52\textwidth}
    \centering
    \textbf{Valley-mixed excitons (bright)}
    \begin{equation}
    \begin{aligned}
        \boldsymbol{D}^{X}_{B+, \boldsymbol{Q}} &= \frac{D^X_B}{\sqrt{2}} \left(\cos\phi_{\boldsymbol{Q}}\,\hat{\boldsymbol{x}} + \sin\phi_{\boldsymbol{Q}}\,\hat{\boldsymbol{y}}\right), \\
        \boldsymbol{D}^{X}_{B-, \boldsymbol{Q}} &= \frac{iD^X_B}{\sqrt{2}} \left(-\sin\phi_{\boldsymbol{Q}}\,\hat{\boldsymbol{x}} + \cos\phi_{\boldsymbol{Q}}\,\hat{\boldsymbol{y}}\right)
    \end{aligned}
    \label{transitiondipole}
    \end{equation}
    
    \textbf{Gray excitons}
    \begin{equation}
        \boldsymbol{D}^{X}_{G, \boldsymbol{Q}} = D^X_G\,\hat{\boldsymbol{z}}.
    \end{equation}
\end{minipage}
\end{figure}

\noindent
Here, the subscript symbols $K$, $K'$, $B\pm$, and $G$ denote the $K$ valley, $K'$ valley, upper/ lower band, and gray exciton, respectively. The magnitude of the dipole moment of BXs (GX) is $D^X_{B}$ ($D^X_G$)~\cite{peng2022tailoring, sanchez2024enhanced}. From Eq.~(\ref{transitionrate}), it becomes clear why a single LG beam, whose TAM is encoded purely in its phase, does not produce a transition rate that reflects SOI or TAM. This is because the transition rate depends on the squared magnitude of the vector potential, which removes all phase information.
%
%
%
%
%
%
\section{Manifestation of SOI in Excitonic Optical Transitions}
\subsection{Optical Transition Pattern of Excitons}
We now can utilize the formulation above to analyze the SOI in excitonic optical transition for both valley-mixed BXs and GX under VVBs. The transition rates can be evaluated with Eq.~(\ref{VVBtrans}), Eq.~(\ref{VVBlongi}), and Eq.~(\ref{transitionrate}), yielding:
\begin{equation}
\begin{aligned}
\left|\tilde{M}^{\sigma, \ell, \sigma', \ell'}_{B+, \boldsymbol{Q}}\right|^2
&=\left( \frac{E_g D^X_B}{2\sqrt{2}\hbar} \right)^2
 \left| \cos(\beta/2)\tilde{F}_{\ell, p}(Q) + e^{-i(\Delta J \phi_{\boldsymbol{Q}} + \alpha)} \sin(\beta/2)\tilde{F}_{\ell', p}(Q) \right|^2, \\
\left|\tilde{M}^{\sigma, \ell, \sigma', \ell'}_{B-, \boldsymbol{Q}}\right|^2
&=\left( \frac{E_g D^X_B}{2\sqrt{2}\hbar} \right)^2
\left| \sigma\cos(\beta/2)\tilde{F}_{\ell, p}(Q) + e^{-i(\Delta J \phi_{\boldsymbol{Q}} + \alpha)} \sigma'\sin(\beta/2)\tilde{F}_{\ell', p}(Q) \right|^2, \\
\left|\tilde{M}^{\sigma, \ell, \sigma', \ell'}_{G, \boldsymbol{Q}}\right|^2
&=\left( \frac{E_g D^X_G}{2\sqrt{2}\hbar} \right)^2
\sin^2 \theta_{\boldsymbol{Q}} \left| \cos(\beta/2)\tilde{F}_{\ell, p}(Q) + e^{-i(\Delta J \phi_{\boldsymbol{Q}} + \alpha)} \sin(\beta/2)\tilde{F}_{\ell', p}(Q) \right|^2,
\end{aligned}
\label{eq:transition_rates_complex}
\end{equation}
In this form, we observe that the SOI term $\Delta J = (\ell' + \sigma') - (\ell + \sigma)$ acts as an additional relative phase modulator with $\phi_{\boldsymbol{Q}}$ dependence. Comparing the cases of valley-mixed BXs and GXs, Eq.~(\ref{firstmatrixelement}) shows that the spin index $\sigma$ can enters the phase of the optical matrix element. Consequently, the $\phi_{\boldsymbol{Q}}$ dependence of the transition dipoles for valley-mixed BXs in Eq.~(\ref{transitiondipole}) provides a mechanism for the emergence of $\Delta J$ under excitation by a VVB state, whose polarization texture also depends on $\phi_{\boldsymbol{Q}}$. As a result, valley-mixed BXs can exhibit interference features analogous to those of GXs.

In contrast, for GXs, the SOI manifests directly through the interference term in Eq.~(\ref{VVBlongi}). This originates from the fact that the transition dipole moment of GXs is oriented along the out-of-plane ($z$) direction and remains isotropic in the in-plane coordinates, making the interference the dominant mechanism responsible for revealing SOI with azimuthal structure.
To explicitly reveal the interference effects, we expand the squared matrix elements for valley-mixed BXs as follows:
% Bright excitons
\begin{equation}
\begin{aligned}
\left|\tilde{M}^{\sigma, \ell, \sigma', \ell'}_{B+, \boldsymbol{Q}}\right|^2
&= \left( \frac{E_g D^X_B}{2\sqrt{2}\hbar} \right)^2
\left\{
\cos^2(\beta/2) \left|\tilde{F}_{\ell, p}(Q)\right|^2
+ \sin^2(\beta/2)\left|\tilde{F}_{\ell', p}(Q)\right|^2 \right. \\
&\hspace{10em}\left.
+ \sin\beta \,
\mathrm{Re}\left[
\tilde{F}_{\ell, p}(Q)\tilde{F}^*_{\ell', p}(Q)
e^{-i(\Delta J \phi_{\boldsymbol{Q}} + \alpha)}
\right]
\right\}, \\
\left|\tilde{M}^{\sigma, \ell, \sigma', \ell'}_{B-, \boldsymbol{Q}}\right|^2
&= \left( \frac{E_g D^X_B}{2\sqrt{2}\hbar} \right)^2
\left\{
\cos^2(\beta/2) \left|\tilde{F}_{\ell, p}(Q)\right|^2
+ \sin^2(\beta/2)\left|\tilde{F}_{\ell', p}(Q)\right|^2 \right. \\
&\hspace{10em}\left.
+ \sigma\sigma'\sin\beta \,
\mathrm{Re}\left[
\tilde{F}_{\ell, p}(Q)\tilde{F}^*_{\ell', p}(Q)
e^{-i(\Delta J \phi_{\boldsymbol{Q}} + \alpha)}
\right]
\right\}.
\end{aligned}
\label{expanded_B}
\end{equation}
Similarly, for GXs, we obtain:
\begin{equation}
\begin{aligned}
\left|\tilde{M}^{\sigma, \ell, \sigma', \ell'}_{G, \boldsymbol{Q}}\right|^2
&= \left( \frac{E_g D^X_G}{2\sqrt{2}\hbar} \right)^2
\sin^2\theta_{\boldsymbol{Q}}
\left\{
\cos^2(\beta/2) \left|\tilde{F}_{\ell, p}(Q)\right|^2
+ \sin^2(\beta/2)\left|\tilde{F}_{\ell', p}(Q)\right|^2 \right. \\
&\hspace{10em}\left.
+ \sin\beta \,
\mathrm{Re}\left[
\tilde{F}_{\ell, p}(Q)\tilde{F}^*_{\ell', p}(Q)
e^{-i(\Delta J \phi_{\boldsymbol{Q}} + \alpha)}
\right]
\right\}.
\end{aligned}
\label{expanded_G}
\end{equation}

In general, the transition rates for BXs and GXs exhibit interference structures analogous to the familiar interference term of the VVB state in Eq.~(\ref{VVBlongi}). In addition, the transition rates of BXs take the same functional form when $\sigma=\sigma'$. Below, we discuss valley-mixed BXs and GX separately, since the physical origins of the SOI term are different. This distinction also affects how the manipulation of the transition rate should be interpreted. Note that one cannot simply set $\sigma=0$ to represent linearly polarized light in this formulation, because the exponential form has been taken to simplify the derivation of the final bracket in Eq.~(\ref{transitionrate}).

\paragraph*{Valley-mixed Bright Excitons : Polarization-Matching}~\\

For valley-mixed BXs, Eq.~(\ref{expanded_B}) shows that using opposite circular polarizations ($\sigma'=-\sigma$) leads to opposite interference behavior for the two valley-mixed branches. In addition, the SOI term $\Delta J$, which does not explicitly appear in the transverse intensity profile of the VVB itself, emerges once the phase factor generated by the dot product between the polarization unit vector and the transition dipole moment is included in Eq.~(\ref{firstmatrixelement}). 
 
Motivated by the polarization texture of the VVB state $|1,-1,-1,1;\alpha,\beta\rangle$, shown in Figure~\ref{fig:VVBs_lm1}, the condition $\Delta J=0$ can be interpreted as a polarization-matching mechanism, illustrated in Figure~\ref{fig:dipole}. In particular, the polarization of the states
\[
|1,-1,-1,1;\alpha=0,\beta=\pi/2\rangle
\quad \text{and} \quad
|1,-1,-1,1;\alpha=\pi,\beta=\pi/2\rangle
\]
matches the directions of the exciton transition dipoles $D^X_B$ and $D^X_B$, respectively (see Figure~\ref{fig:dipole}). This enables selective excitation of the upper and lower exciton bands with the relative phase $\alpha$ as shown in Figure~\ref{fig:BXtransition}.
\begin{figure}[H]
        \centering
        \includegraphics[width=0.65\linewidth]{dipole.pdf}
        \caption{Illustration of the transition dipole moments of valley-mixed BXs~\cite{peng2022tailoring}, and their matching with the transverse polarization of the VVB states $|1,-1,-1,1;\alpha=0,\beta=\pi/2\rangle$ and $|1,-1,-1,1;\alpha=\pi,\beta=\pi/2\rangle$ in momentum space.}
        \label{fig:dipole}
\end{figure}
\begin{figure}[H]
        \centering
        \includegraphics[width=0.75\linewidth]{BXtransition.pdf}
        \caption{Normalized transition rate of valley-mixed BXs excited by the VVB state $|1,-1,-1,1;\alpha,\beta\rangle$ in momentum space. The scale bar indicates the considered $0.1$ times light cone range for WSe$_2$-monolayers, $Q_c = 8.6~\mu\mathrm{m}^{-1}$}
        \label{fig:BXtransition}
\end{figure}
However, since the valley splitting between valley-mixed BXs is only a few meV and is typically regarded as spectrally unresolvable~\cite{schneider2021optical}, direct experimental verification of selective excitation may be difficult. Therefore, we consider the transition rate for the total BXs under the condition $(\sigma'=-\sigma)$, which is given by
\begin{equation}
\begin{aligned}
\Gamma^{\sigma,\ell,-\sigma,\ell'}_{B\pm,\boldsymbol{Q}}(\alpha,\beta)
&\propto
\sum_{S=B\pm}
\left|
\tilde{M}^{\sigma,\ell,-\sigma,\ell'}_{S,\boldsymbol{Q}}
\right|^2 \\
&=
\left(
\frac{E_g D_B^X}{2\hbar}
\right)^2
\left[
\cos^2\!\left(\frac{\beta}{2}\right)
\left|\tilde{F}_{\ell,p}(Q)\right|^2
+
\sin^2\!\left(\frac{\beta}{2}\right)
\left|\tilde{F}_{\ell',p}(Q)\right|^2
\right].
\end{aligned}
\label{transitionrate_bothBX}
\end{equation}
The formulation shows the cancellation of the two interference terms, leading to the disappearance of both the $\alpha$ dependence and the $\Delta J$ dependence. The resulting distribution is shown in Figure~\ref{fig:BXtransition(sum)}, which exhibits an identical donut-shaped profile over the Poincar\'e sphere.
\begin{figure}[H]
    \centering
    \includegraphics[width=0.75\linewidth]{BXtransition(sum).pdf}
    \caption{Normalized transition rate of total BXs excited by the VVB state $|1,-1,-1,1;\alpha,\beta\rangle$ in momentum space.}
    \label{fig:BXtransition(sum)}
\end{figure}

It should be noted that the condition $\Delta J=0$ is not unique to this case. It can also be satisfied when $\ell'=\ell$ together with $\sigma'=\sigma$. However, in that situation the polarization-matching mechanism is absent. Instead, selective excitation is achieved simply by tuning the relative superposition parameters, effectively switching the excitation channel on and off, as shown in Figure~\ref{fig:dipole_same} and Figure~\ref{fig:BXtransition_same}. Consequently, the excitation behavior of the two BX branches becomes identical, consistent with Eq.~(\ref{expanded_B}).

\begin{figure}[H]
    \centering
    \begin{minipage}{0.39\textwidth}
        \centering
        \includegraphics[width=\linewidth]{dipole(same).pdf}
        \caption{Distribution of the transverse component for the VVB state $|1,1,1,1;\alpha,\beta\rangle$ in momentum space.}
        \label{fig:dipole_same}
    \end{minipage}
    \hfill
    \begin{minipage}{0.57\textwidth}
        \centering
        \includegraphics[width=\linewidth]{BXtransition(same).pdf}
        \caption{Normalized transition rate of valley-mixed BXs excited by the VVB state $|1,1,1,1;\alpha,\beta\rangle$ in momentum space.}
        \label{fig:BXtransition_same}
    \end{minipage}
\end{figure}

\paragraph*{Gray Excitons: Manifestation of SOI}~\\

For gray excitons, whose transition dipole moments possess only a $z$ component, no additional phase factor is generated by the dot product between the polarization vector and the transition dipole moment in Eq.~(\ref{firstmatrixelement}). Therefore, the SOI term $\Delta J$ in Eq.~(\ref{expanded_G}) originates entirely from the longitudinal component of the VVB state in Eq.~(\ref{VVBlongi}). As a result, GX provide an ideal probe of the SOI of light, since their transition rates directly reflect the intrinsic angular-momentum structure of the optical field. At the same time, they offer a possible route for exciting GX using paraxial structured beams.

Further exploit the intensity texture of the VVB state $|1,-1,-1,1;\alpha,\beta\rangle$, shown in Figure~\ref{fig:VVBs_lm1}. In particular, complete switching of the excitation on and off can be realized by tuning the relative phase between $\alpha=0$ and $\alpha=\pi$ under excitation by the state
\[
|1,-1,-1,1;\alpha,\beta=\pi/2\rangle,
\]
as shown in Figure~\ref{fig:GXtransition}.
\begin{figure}[H]
    \centering
    \includegraphics[width=0.75\linewidth]{GXtransition.pdf}
    \caption{Normalized transition rate of GXs excited by the VVB state $|1,-1,-1,1;\alpha,\beta\rangle$ in momentum space.}
    \label{fig:GXtransition}
\end{figure}

Beyond complete switching, one may also use the case $\Delta J=1$ to excite gray excitons with selected finite center-of-mass momentum. For example, consider the state
\[
|1,-1,-1,2;\alpha,\beta=\pi/2\rangle,
\]
for which
\[
\Delta J=(\ell'+\sigma')-(\ell+\sigma)=1.
\]

First, the radial factor $\tilde{F}_{\ell p}(q_\parallel)$ of the two constituent LG beams are no longer identical, so the transition rates at $\beta=0$ and $\beta=\pi$ (the north and south poles of the Poincar\'e sphere) is different. Figure~\ref{fig:GXtransition(L+3)} shows that the LG beam with larger $\ell$ produces a larger-radius donut-shaped transition-rate profile in the momentum space.
\begin{figure}[H]
    \centering
    \begin{minipage}{0.37\textwidth}
        \centering
        % trim=左 下 右 上, clip
        \includegraphics[width=\linewidth, trim=4cm 0cm 4cm 0cm, clip]{GXtransition(L+3).pdf}
        \caption{Normalized transition rate of GXs excited by the VVB state $|1,-1,-1,2;\alpha,\beta\rangle$ in momentum space at the north/south poles of the Poincar\'e sphere.}
        \label{fig:GXtransition(L+3)}
    \end{minipage}
    \hfill
    \begin{minipage}{0.59\textwidth}
        \centering
        \includegraphics[width=\linewidth]{GXtransition(L+3)_2.pdf}
        \caption{Normalized transition rate of GXs excited by the VVB state $|1,-1,-1,2;\alpha,\beta\rangle$ in momentum space along the equator of the Poincar\'e sphere. The scale and color bars are the same as in Figure~\ref{fig:GXtransition(L+3)}.}
        \label{fig:GXtransition(L+3)_2}
    \end{minipage}
\end{figure}
Second, although the two LG beams do not fully overlap, the condition $\Delta J=1$ still induces one full interference cycle in the azimuthal direction. Consequently, by tuning $\alpha$, one can selectively excite GX within a finite range of center-of-mass momentum at $\beta=\pi/2$ (the equator of the Poincar\'e sphere), as shown in Figure~\ref{fig:GXtransition(L+3)_2}.
%
%
%
%
%
%
\subsection{Selective Excitation of Excitons}

By summing the contributions in momentum space excited by the VVBs, we obtain the total transition rate, which represents the intensity of the material absorption. Under the paraxial approximation, where $\theta_{\boldsymbol{Q}} = \frac{Q}{\sqrt{Q^2+q_0^2}} \approx \frac{Q}{q_0}$ is taken to simplify the integral, the total transition rate can be evaluated as:
\[
\Gamma_{S}^{\sigma, \ell, \sigma', \ell'} 
\equiv 
\sum_{\boldsymbol{Q}} 
\Gamma_{S, \boldsymbol{Q}}^{\sigma, \ell, \sigma', \ell'} 
=
\frac{\Omega}{(2\pi)^2} 
\int_0^\infty \int_0^{2\pi} 
\Gamma_{S, \boldsymbol{Q}}^{\sigma, \ell, \sigma', \ell'} 
\, Q \, \mathrm{d}\phi_{\boldsymbol{Q}} \, \mathrm{d}Q.
\]
The total transition rates for the upper-band BXs, lower-band BXs, and GXs are obtained as
\begin{equation}
\begin{aligned}
\Gamma^{\sigma,\ell,\sigma',\ell'}_{B+}(\alpha,\beta)
&\propto \left|\tilde{M}^{\sigma,\ell,\sigma',\ell'}_{B+}\right|^2 \\
&= \frac{\Phi}{8} \left( \frac{E_g D^X_B}{a_0\Omega\hbar} \right)^2 
\left\{1+\delta_{\Delta J, 0}\sin\beta\cos\left[\frac{\pi}{2}(|\ell'|-|\ell|)+\alpha\right]\frac{\left(\frac{|\ell'|+|\ell|}{2}\right)!}{\sqrt{|\ell'|!|\ell|!}}\right\} , \\
\Gamma^{\sigma,\ell,\sigma',\ell'}_{B-}(\alpha,\beta) 
&\propto \left|\tilde{M}^{\sigma,\ell,\sigma',\ell'}_{B-}\right|^2 \\
&=\frac{\Phi}{8} \left( \frac{E_g D^X_B}{a_0\Omega\hbar} \right)^2 
\left\{1+\delta_{\Delta J, 0}\sigma\sigma'\sin\beta\cos\left[\frac{\pi}{2}(|\ell'|-|\ell|)+\alpha\right]\frac{\left(\frac{|\ell'|+|\ell|}{2}\right)!}{\sqrt{|\ell'|!|\ell|!}}\right\} , \\
\Gamma^{\sigma,\ell,\sigma',\ell'}_{G}(\alpha,\beta) 
&\propto \left|\tilde{M}^{\sigma,\ell,\sigma',\ell'}_{G}\right|^2 \\
&=\frac{\Phi}{4} \left( \frac{E_g D^X_G}{a_0 q_0 W_0\Omega\hbar} \right)^2 
\left\{\left(1+\cos^2\frac{\beta}{2}|\ell|+\sin^2\frac{\beta}{2}|\ell'|\right) \right. \\
&\hspace{7em} \left. + \delta_{\Delta J, 0}\sin\beta\cos\left[\frac{\pi}{2}(|\ell'|-|\ell|) +\alpha\right]\frac{\left(\frac{|\ell'|+|\ell|}{2}+1\right)!}{\sqrt{|\ell'|!|\ell|!}}\right\} ,
\end{aligned}
\label{total_transition_rate}
\end{equation}
where the constant radiant power is $\Phi = (a_0 A_0^{\mathrm{LG}} W_0)^2$, and $a_0 = \frac{\omega}{\sqrt{2c\mu_0}}$.

At first glance, the SOI appears to be absent, since $\ell$ and $\sigma$ are not explicitly coupled in the second term of all three expressions. However, the trigonometric factor can be reinterpreted as
\begin{equation}
\sin\beta\cos\left[\frac{\pi}{2}(|\ell'|-|\ell|)+\alpha\right]
=
\begin{pmatrix}
{\rm Re}\left[e^{i\left(|\ell'|-|\ell|\right)}\right], 
{\rm Im}\left[e^{i\left(|\ell'|-|\ell|\right)}\right],
0
\end{pmatrix}
\cdot
\begin{pmatrix}
\sin\beta\cos\alpha\\
\sin\beta\sin\alpha\\
\cos\beta
\end{pmatrix},
\end{equation}
which reveals its origin as a phase correction between the complex components of two VVB states located on the poles of Poincar\'e sphere.

More specifically, the complex phase of the radial factor $\tilde{F}_{\ell,p}(Q)$ takes the form $(-i)^{|\ell|}$. As a result, the interference is governed by the difference in $|\ell|$, which introduces a phase shift across whole momentum space, together with the VVB state specified by $(\alpha,\beta)$.

At the same time, the contribution from SOI (or, for BXs, the polarization-matching mechanism) is encoded through $\delta_{\Delta J, 0}$. When $\Delta J \neq 0$, the resulting azimuthal phase mismatch—such as that illustrated in Figure~\ref{fig:GXtransition(L+3)_2}—leads to destructive averaging over momentum space, effectively given no net contribution to the total transition rate.

From this perspective, the factorial term in Eq.~(\ref{total_transition_rate}) can be interpreted as a geometric correction arising from the mismatch between the radial factors $\tilde{F}_{\ell,p}(Q)$ of the two LG beams. In particular, when $|\ell| = |\ell'|$, it reduces to simple constants, yielding $1$ for BXs and $1 + |\ell|$ for GXs. These values are consistent with the corresponding prefactors in the non-interference contribution (i.e., the term independent of $\delta_{\Delta J, 0}$).

As for the beam waist $W_0$, which appears in the denominator of the prefactor for the GX total transition rate, it reflects that smaller $W_0$ enhances the longitudinal field component and thus increases the excitation efficiency of GXs.

On the other hand, the total transition rate for the total BXs under the condition $(\sigma'=-\sigma)$ can be similarly written as
\begin{equation}
\Gamma^{\sigma,\ell,-\sigma,\ell'}_{B\pm}(\alpha,\beta)
\propto \sum_{S=B\pm} \left|\tilde{M}^{\sigma,\ell,-\sigma,\ell'}_{S}\right|^2 
= \frac{\Phi}{4} \left( \frac{E_g D^X_B}{a_0\Omega\hbar} \right)^2 .
\label{totaltransitionrate_bothBX}
\end{equation}

Therefore, one can exploit both Eq.~(\ref{total_transition_rate}) and Eq.~(\ref{totaltransitionrate_bothBX}) to analyze the optical transition rates. In particular, we focus on the equator of the Poincar\'e sphere ($\beta=\pi/2$), where the VVB corresponds to a maximal superposition of two LG beams. By varying the azimuthal angle $\alpha$ along the equator, we examine the dependence of the total transition rate on $\ell$ (with $|\ell|=|\ell'|$) using the VVB states $|1,\ell,-1,-\ell;\alpha,\beta=\pi/2\rangle$, as shown in Figure~\ref{fig:totaltransition}.
\begin{figure}[H]
    \centering
    \includegraphics[width=1.0\linewidth]{totaltransition.pdf}
    \caption{Normalized total transition rates of BXs, total BXs, and GX excited by a VVB state $|1,\ell,-1,-\ell;\alpha,\beta=\pi/2\rangle$, for $\ell \in \mathbb{Z}$ with $\ell \in [-5,5]$, and for azimuthal angles $\alpha = 0, \pi/2, \pi$.}
    \label{fig:totaltransition}
\end{figure}

In Figure~\ref{fig:totaltransition}, the total transition rate for each exciton exhibits a clear dependence on $\alpha$ when the condition $\Delta J = 0$ is satisfied at $\ell=-1$. As shown within the same column, varying $\alpha$ from $0$ to $\pi/2$ to $\pi$ leads to modulation of the transition rate. This behavior is consistent with the results in Figure~\ref{fig:BXtransition} and Figure~\ref{fig:GXtransition}, where the transition rate exhibits fully constructive or destructive interference over momentum space at $\alpha = 0$ or $\alpha = \pi$.

In contrast, the total transition rate of total BXs becomes independent of $\alpha$ and $\ell$. This is due to the cancellation between the contributions of the upper and lower BX bands, as reflected in Figure~\ref{fig:BXtransition_same}.

Apart from this special cancellation, BXs generally exhibit independence from $\ell$, unlike GXs. This difference originates from the additional $\boldsymbol{Q}$-dependent prefactor in the GX case, specifically $\sin\theta_{\boldsymbol{Q}} \approx Q/q_0$ in Eq.~(\ref{expanded_G}). Consequently, Eq.~(\ref{total_transition_rate}) shows a dependence on $\ell$ even in the absence of interference effects.

Motivated by the special properties of the VVB state at the equator of the Poincar\'e sphere, $|1,-1,-1,1;\alpha,\beta=\pi/2\rangle$, one can achieve smooth and selective excitation over half of the equator ($\alpha = 0$ to $\pi$), as shown in Figure~\ref{fig:totaltransition_special}.

\begin{figure}[H]
    \centering
    \includegraphics[width=0.8\linewidth]{totaltransition_special.pdf}
    \caption{Normalized total transition rates of BXs, total BXs, and GXs excited by a VVB state $|1,-1,-1,1;\alpha,\beta=\pi/2\rangle$ for azimuthal angles $\alpha = 0$ to $2\pi$.}
    \label{fig:totaltransition_special}
\end{figure}

These observations originate from SOI associated with the longitudinal component of the LG beam, together with a polarization-matching mechanism. The former arises from the gauge transformation, specifically the second term in Eq.~(\ref{LCrelation2}), as well as the intrinsic structure of the LG beam. The latter originates from the correspondence between the transition dipole structure of the upper and lower BX bands and the polarization textures of the VVB state. In both cases, Fermi's golden rule provides a theoretical framework for interpreting these results and establishes a basis for further studies involving different VVB configurations, excitonic species, and other comprehensive investigations.






